\chapter{Monitoring and Control}

Per evitare che problemi come ritardi o costi aggiuntivi compromettessero il progetto, MaraffaOnline è stato monitorato e controllato in modo continuo lungo tutti i sette mesi. Il monitoraggio si è articolato su tre livelli: quotidiano, per i sottosistemi Agile, tramite Daily Standup di 15 minuti; settimanale, tramite Project Status Meeting di un'ora con lo sponsor; e per sprint, tramite Sprint Review e Retrospective al termine di ogni iterazione.

PlayHeritage Labs ha promosso un clima di apertura in cui i problemi vengono segnalati con tempestività, così da individuarli prima che si accumulino. A supporto, ogni sviluppatore aggiorna quotidianamente lo stato dei propri task sulla board di project management, e un Issue Log su Notion tiene traccia di tutti i problemi riscontrati durante il progetto.

\section{Reporting}

Per comunicare lo stato di avanzamento allo sponsor si è adottato il formato \textbf{Stoplight Report}: a ciascuna area critica (Scope, Schedule, Budget, Quality, Risks) viene assegnato un colore --- verde (in linea), giallo (scostamento sotto controllo), rosso (intervento necessario) --- con una nota che ne motiva la valutazione. I colori gialli segnalano le aree da monitorare prima che diventino rosse.

Il controllo di costi e tempi è stato affidato all'\textbf{Earned Value Management}, che confronta il valore pianificato (PV), il valore realizzato (EV) e il costo effettivo (AC), derivandone gli indici di efficienza CPI e SPI. Al termine del progetto il consumo è stato di €22.750 su €25.000, con CPI e SPI prossimi a 1: segno di una buona aderenza tra pianificato ed effettivo. L'analisi Earned Value mensile completa, gli Stoplight Report settimanali e le metriche di velocity e di qualità (test coverage stabilmente sopra l'85\%) sono riportati nel documento di dettaglio Monitoring \& Control allegato alla relazione.

\section{Problem Solving e Gestione dei Rischi}

Di fronte a uno scostamento, la strategia di intervento ha seguito una scala graduata: prima si valuta se lo slack disponibile assorbe il problema, poi si esaminano le dipendenze tra i task, e solo se necessario si rinegozia con lo sponsor una rischedulazione. L'esempio principale è stato il ritardo di 3 giorni accumulato sul Frontend Dashboard e Creazione Stanza a Gennaio 2026 --- attività sul percorso critico, a monte del Tavolo da Gioco --- riassorbito entro Febbraio grazie a un recovery plan (pair programming) senza impatto sulla milestone finale.

I rischi identificati nella Risk Rating Matrix (\allegatolink{Documentazione/Scoping/Allegato2.4-RiskAnalysis.pdf}{Allegato 2.4}) sono stati monitorati a ogni Project Status Meeting. Il rischio più critico --- l'esperienza limitata del team con le tecnologie WebSocket --- è stato affrontato con uno spike tecnico e un proof of concept in Sprint 2 (latenza misurata a 180ms): la sua probabilità si è ridotta portando il rating da 16 a 8, fino a chiuderlo come mitigato dopo il load test dello Sprint 6.

Grazie a questo monitoraggio multilivello tutti i criteri definiti nelle Conditions of Satisfaction sono stati rispettati: la soluzione è stata collaudata internamente, ha superato la fase di UAT ed è giunta pronta al lancio pubblico del 15 Maggio 2026. La chiusura formale del progetto è descritta nel capitolo successivo.
